%% Chapter 10 — Charts and Figures Catalogue
\chapter{Charts and Figures Catalogue}
\label{ch:figures}

\section{Introduction}

This chapter provides a consolidated catalogue of all charts, maps, and figures generated by the \BDE{} pipeline. Each figure is presented with an extended caption describing the data source, generation method, and key interpretive notes. Figures referenced in earlier chapters are reproduced here in higher resolution with additional annotation.


\section{Site and Context Maps}

\subsection{Site Location Overview}

\begin{figure}[H]
\centering
\begin{tikzpicture}
  \draw[dreamlab-light,fill=dreamlab-light,rounded corners] (0,0) rectangle (14,9);
  \node[dreamlab-grey,font=\sffamily\large] at (7,4.5) {../output/figures/site\_location\_overview\_hires.png};
\end{tikzpicture}
\caption[Site location overview map.]{%
\textbf{Site location overview.} The study area is centred on Pingle, Taylor's Lane, Ashleyhay (SK~310~525), shown within the wider Derbyshire landscape. The 1\,km radius study boundary is indicated by the red circle. Inset: regional context showing the study area's position relative to the Peak District National Park boundary and the Derwent Valley Mills World Heritage Site. Base mapping: OS MasterMap Topography Layer, \copyright\ Crown Copyright 2026.%
}
\label{fig:cat-site-overview}
\end{figure}

\subsection{Environmental Designations Map}

\begin{figure}[H]
\centering
\begin{tikzpicture}
  \draw[dreamlab-light,fill=dreamlab-light,rounded corners] (0,0) rectangle (14,9);
  \node[dreamlab-grey,font=\sffamily\large] at (7,4.5) {../output/figures/designations\_map.png};
\end{tikzpicture}
\caption[Environmental designations within and proximate to the study area.]{%
\textbf{Environmental designations.} Statutory designations (SSSI — hatched red), non-statutory designations (LWS — hatched green), Ancient Woodland Inventory parcels (brown), and the Derwent Valley Mills WHS buffer zone (purple dashed) are overlaid on the study area boundary. Data: Natural England Open Data, accessed March 2026.%
}
\label{fig:cat-designations}
\end{figure}


\section{Remote Sensing Composites}

\subsection{True-Colour Composites}

\begin{figure}[H]
\centering
\begin{tikzpicture}
  \draw[dreamlab-light,fill=dreamlab-light,rounded corners] (0,0) rectangle (14,6);
  \node[dreamlab-grey,font=\sffamily] at (3.5,3) {2016 RGB Composite};
  \draw[dreamlab-grey,dashed] (7,0) -- (7,6);
  \node[dreamlab-grey,font=\sffamily] at (10.5,3) {2026 RGB Composite};
\end{tikzpicture}
\caption[Side-by-side true-colour composites 2016 vs 2026.]{%
\textbf{Growing-season true-colour composites.} Sentinel-2 B4-B3-B2 median composites for the 2016 (left) and 2026 (right) growing seasons. Composites generated from \dataplaceholder{N 2016} and \dataplaceholder{N 2026} cloud-free scenes respectively. Histogram-matched for visual comparability.%
}
\label{fig:cat-rgb-comparison}
\end{figure}

\subsection{False-Colour Composites}

\begin{figure}[H]
\centering
\begin{tikzpicture}
  \draw[dreamlab-light,fill=dreamlab-light,rounded corners] (0,0) rectangle (14,6);
  \node[dreamlab-grey,font=\sffamily] at (3.5,3) {2016 NIR-R-G};
  \draw[dreamlab-grey,dashed] (7,0) -- (7,6);
  \node[dreamlab-grey,font=\sffamily] at (10.5,3) {2026 NIR-R-G};
\end{tikzpicture}
\caption[False-colour composites highlighting vegetation vigour.]{%
\textbf{False-colour composites (NIR-Red-Green).} Healthy, dense vegetation appears bright red; bare soil and built surfaces appear grey-blue. This band combination enhances the distinction between habitat types and highlights changes in vegetation density between epochs.%
}
\label{fig:cat-fc-comparison}
\end{figure}

\subsection{NDVI Composites}

\begin{figure}[H]
\centering
\begin{tikzpicture}
  \draw[dreamlab-light,fill=dreamlab-light,rounded corners] (0,0) rectangle (14,6);
  \node[dreamlab-grey,font=\sffamily] at (3.5,3) {2016 NDVI};
  \draw[dreamlab-grey,dashed] (7,0) -- (7,6);
  \node[dreamlab-grey,font=\sffamily] at (10.5,3) {2026 NDVI};
\end{tikzpicture}
\caption[NDVI composites 2016 vs 2026.]{%
\textbf{Normalised Difference Vegetation Index composites.} NDVI values range from $-$1 (water, bare) to $+$1 (dense vegetation). The green-to-brown colour ramp highlights spatial variation in vegetation productivity. Areas of NDVI increase between epochs may indicate habitat condition improvement; areas of decrease may indicate degradation or conversion.%
}
\label{fig:cat-ndvi-comparison}
\end{figure}


\section{Classification Maps}

\subsection{Habitat Classification Maps}

\begin{figure}[H]
\centering
\begin{tikzpicture}
  \draw[dreamlab-light,fill=dreamlab-light,rounded corners] (0,0) rectangle (14,9);
  \node[dreamlab-grey,font=\sffamily\large] at (7,4.5) {../output/figures/habitat\_classification\_paired.png};
\end{tikzpicture}
\caption[Paired habitat classification maps 2016 and 2026.]{%
\textbf{UKHab Level~3 classification maps.} Left: 2016 baseline (\Tzero). Right: 2026 present (\Tone). Colour coding follows the UKHab standard palette. Parcels with changed classification are highlighted with bold outlines on the 2026 map. Legend shows all habitat classes present in either epoch.%
}
\label{fig:cat-classification-paired}
\end{figure}

\subsection{Classification Confidence Map}

\begin{figure}[H]
\centering
\begin{tikzpicture}
  \draw[dreamlab-light,fill=dreamlab-light,rounded corners] (0,0) rectangle (14,9);
  \node[dreamlab-grey,font=\sffamily\large] at (7,4.5) {../output/figures/classification\_confidence\_map.png};
\end{tikzpicture}
\caption[Random Forest classification confidence.]{%
\textbf{Classification confidence.} Per-parcel maximum class probability from the Random Forest classifier. Values range from 0.5 (minimum for a binary choice) to 1.0 (certain). Parcels with confidence below 0.7 are highlighted with hatching and warrant particular scrutiny in the delta analysis.%
}
\label{fig:cat-confidence-map}
\end{figure}


\section{Biodiversity Unit Maps}

\subsection{BU Choropleth Maps}

\begin{figure}[H]
\centering
\begin{tikzpicture}
  \draw[dreamlab-light,fill=dreamlab-light,rounded corners] (0,0) rectangle (14,9);
  \node[dreamlab-grey,font=\sffamily\large] at (7,4.5) {../output/figures/bu\_choropleth\_paired.png};
\end{tikzpicture}
\caption[Paired BU choropleth maps.]{%
\textbf{Biodiversity unit density maps.} BU per hectare displayed as a graduated colour map: light yellow (low BU/ha) through orange to dark red (high BU/ha). Left: 2016. Right: 2026. The colour scale is identical in both maps for direct visual comparison.%
}
\label{fig:cat-bu-choropleth}
\end{figure}

\subsection{Delta Map}

\begin{figure}[H]
\centering
\begin{tikzpicture}
  \draw[dreamlab-light,fill=dreamlab-light,rounded corners] (0,0) rectangle (14,9);
  \node[dreamlab-grey,font=\sffamily\large] at (7,4.5) {../output/figures/delta\_map\_annotated.png};
\end{tikzpicture}
\caption[Annotated parcel-level delta map.]{%
\textbf{Parcel-level $\Delta$BU map (annotated).} Diverging colour scale: green indicates biodiversity gain, white indicates zero change, red indicates loss. Parcels with $|\Delta\text{BU}| > $ \dataplaceholder{THRESHOLD} are labelled with their delta value. Hatched parcels indicate habitat class change between epochs. The five parcels with the largest absolute delta are labelled A--E and discussed in the worked examples (Chapter~\ref{ch:delta}).%
}
\label{fig:cat-delta-annotated}
\end{figure}


\section{Statistical Charts}

\subsection{Classification Confusion Matrix}

\begin{figure}[H]
\centering
\begin{tikzpicture}
  \draw[dreamlab-light,fill=dreamlab-light,rounded corners] (0,0) rectangle (12,10);
  \node[dreamlab-grey,font=\sffamily\large] at (6,5) {../output/figures/confusion\_matrix\_2016.png};
\end{tikzpicture}
\caption[Classification confusion matrix (2016 epoch).]{%
\textbf{Confusion matrix for the 2016 habitat classification.} Rows: true class (validation labels). Columns: predicted class. Cell values are normalised by row (recall). Diagonal values represent per-class recall; off-diagonal values represent misclassification rates. The most common confusion is between modified grassland (g4) and other neutral grassland (g3c).%
}
\label{fig:cat-confusion-2016}
\end{figure}

\subsection{Monte Carlo Distribution}

\begin{figure}[H]
\centering
\begin{tikzpicture}
  \draw[dreamlab-light,fill=dreamlab-light,rounded corners] (0,0) rectangle (12,7);
  \node[dreamlab-grey,font=\sffamily\large] at (6,3.5) {../output/figures/monte\_carlo\_full\_distribution.png};
\end{tikzpicture}
\caption[Monte Carlo delta distribution with annotations.]{%
\textbf{Monte Carlo simulation: probability density of total $\Delta$BU.} 10,000 iterations. The nominal estimate (blue dashed line) falls at \dataplaceholder{NOM} BU. The 95\% confidence interval [\dataplaceholder{CI LOW}, \dataplaceholder{CI HIGH}] is shaded. The red vertical line at zero separates gain from loss territory.%
}
\label{fig:cat-mc-full}
\end{figure}

\subsection{Sensitivity Tornado Diagram}

\begin{figure}[H]
\centering
\begin{tikzpicture}
  \draw[dreamlab-light,fill=dreamlab-light,rounded corners] (0,0) rectangle (12,7);
  \node[dreamlab-grey,font=\sffamily\large] at (6,3.5) {../output/figures/tornado\_sensitivity\_full.png};
\end{tikzpicture}
\caption[Sensitivity analysis tornado diagram.]{%
\textbf{One-at-a-time sensitivity analysis.} Each horizontal bar shows the range of total $\Delta$BU when the corresponding parameter is varied from its low to high plausible bound while all other parameters are held at nominal. The condition threshold is the most influential parameter, followed by classification probability cutoff.%
}
\label{fig:cat-tornado}
\end{figure}

\subsection{LISA Cluster Map}

\begin{figure}[H]
\centering
\begin{tikzpicture}
  \draw[dreamlab-light,fill=dreamlab-light,rounded corners] (0,0) rectangle (14,9);
  \node[dreamlab-grey,font=\sffamily\large] at (7,4.5) {../output/figures/lisa\_cluster\_map\_full.png};
\end{tikzpicture}
\caption[LISA spatial autocorrelation cluster map.]{%
\textbf{Local Indicators of Spatial Association (LISA) cluster map.} Red: High-High clusters (biodiversity gain hotspots). Blue: Low-Low clusters (loss coldspots). Light red/blue: spatial outliers. Grey: not significant (p $>$ 0.05). The gain hotspot in the north-western quadrant corresponds to the limestone grassland area showing condition improvement.%
}
\label{fig:cat-lisa}
\end{figure}


\section{Linear Feature Charts}

\begin{figure}[H]
\centering
\begin{tikzpicture}
  \draw[dreamlab-light,fill=dreamlab-light,rounded corners] (0,0) rectangle (12,7);
  \node[dreamlab-grey,font=\sffamily\large] at (6,3.5) {../output/figures/hedgerow\_condition\_comparison.png};
\end{tikzpicture}
\caption[Hedgerow condition comparison 2016 vs 2026.]{%
\textbf{Hedgerow condition distribution comparison.} Stacked bar charts showing the proportion of total hedgerow length in each condition category at \Tzero{} (left) and \Tone{} (right). The shift towards higher condition categories at \Tone{} indicates overall hedgerow condition improvement across the study area.%
}
\label{fig:cat-hedge-condition}
\end{figure}


\section{Cross-Examination Charts}

\begin{figure}[H]
\centering
\begin{tikzpicture}
  \draw[dreamlab-light,fill=dreamlab-light,rounded corners] (0,0) rectangle (12,7);
  \node[dreamlab-grey,font=\sffamily\large] at (6,3.5) {../output/figures/validation\_agreement\_summary.png};
\end{tikzpicture}
\caption[Validation agreement summary across data sources.]{%
\textbf{Cross-examination agreement summary.} Bar chart comparing classification agreement (Level~3 and Level~2) and condition agreement (exact and within-one) across the three validation data sources: professional ecology surveys, University of Derby thesis, and PHI cross-reference.%
}
\label{fig:cat-validation-summary}
\end{figure}
